% !TeX root = ../../金融经济学_王江_讲义及习题.tex
% 知识点：references/金融经济学/金融经济学-王江_261013.pdf，印刷页 125–137（小结续至 138）
% 习题：references/金融经济学/史上最全最完整的金融经济学答案(王江版).pdf，页 38–42（印刷页 37–41）
% Ross 定理必要性展开有限状态、非退化内点情形；9.5—9.7 的外部偏好刻画按教材保留。

\lecture{组合选择（二）}
\label{chap:wj-nine}
\begingroup
\raggedbottom
\renewcommand{\thesection}{\thechapter.\arabic{section}}
\renewcommand{\thesubsection}{\thesection.\arabic{subsection}}
\renewcommand{\thesubsubsection}{\thesubsection.\arabic{subsubsection}}
\ctexset{
 section={number=\thesection,aftername=\quad},
 subsection={number=\thesubsection,aftername=\quad},
 subsubsection={number=\thesubsubsection,aftername=\quad}
}
\setcounter{equation}{0}
\newcounter{wjNineDefinition}
\newcounter{wjNineTheorem}
\newcounter{wjNineExample}
\newcounter{wjNineExercise}
\renewcommand{\thewjNineDefinition}{\thechapter.\arabic{wjNineDefinition}}
\renewcommand{\thewjNineTheorem}{\thechapter.\arabic{wjNineTheorem}}
\renewcommand{\thewjNineExample}{\thechapter.\arabic{wjNineExample}}
\renewcommand{\thewjNineExercise}{\thechapter.\arabic{wjNineExercise}}
\newcommand{\WJNineDefinition}[2]{\refstepcounter{wjNineDefinition}\label{#1}\textbf{定义 \thewjNineDefinition#2}\quad}
\newcommand{\WJNineTheorem}[2]{\refstepcounter{wjNineTheorem}\label{#1}\textbf{定理 \thewjNineTheorem#2}\quad}
\newcommand{\WJNineExample}[2]{\refstepcounter{wjNineExample}\label{#1}\textbf{例 \thewjNineExample#2}\quad}
\newcommand{\WJNineExercise}[2]{\refstepcounter{wjNineExercise}\label{#1}\problemno{\thewjNineExercise#2}}


\ifdefined\WJHandoutKnowledge

\WJRef{chap:wj-eight}{第}[讲（组合选择一）]在一般期望效用下讨论最优组合的性质。本章进一步限制收益率分布或参与者偏好，回答两个问题：什么时候所有某类参与者都对两个组合给出相同排序？什么时候不同参与者只需持有少数几只共同基金？前者引出\textbf{随机占优}，后者引出\textbf{共同基金分离}。

\section{随机占优}\label{sec:wj-nine-1}
\subsection{一阶随机占优}
\subsubsection{定义与等价判据}
记 $\tilde r_A,\tilde r_B$ 为净收益率，$\tilde x_A=1+\tilde r_A$、$\tilde x_B=1+\tilde r_B$ 为总收益率，$F_A,F_B$ 为总收益率的累积分布函数。固定相同初始投资额后，可以把财富尺度吸收到效用函数中，直接比较收益率。

\begin{tcolorbox}[knowledge,breakable=false]
\WJNineDefinition{def:wj-nine-1}{} 如果对每个递增效用函数（$u'\ge0$）都有
\[
 \mathrm E[u(\tilde x_A)]\ge\mathrm E[u(\tilde x_B)],
\]
则称 $A$ 对 $B$ \textbf{一阶随机占优}（FSD），记作 $A\succeq_{\mathrm{FSD}}B$。
\end{tcolorbox}
这里不限制风险厌恶程度：风险厌恶、风险中性乃至风险喜好的参与者，只要喜欢更多收益，都会弱偏好 $A$。比较的是\textbf{全部投资于组合 $A$ 与全部投资于组合 $B$}，而非已经得到了同时交易两者时的最优头寸。

\begin{tcolorbox}[knowledge,breakable=false]
\WJNineTheorem{thm:wj-nine-1}{} 下列叙述等价：
\begin{enumerate}
 \item $A\succeq_{\mathrm{FSD}}B$；
 \item 对所有 $x$，$F_A(x)\le F_B(x)$；
 \item 可以在同一概率空间构造同分布版本，使
 $\tilde x_A\sim_d\tilde x_B+\tilde d$，其中 $\tilde d\ge0$。
\end{enumerate}
\end{tcolorbox}
$\sim_d$ 表示\textbf{分布相同}。第 3 条是关于存在某种配对的陈述，不表示实际市场中同一状态下的 $A$ 必定比 $B$ 好，也不要求 $\tilde d$ 与 $\tilde x_B$ 独立。

% 课堂照片补充：用同一阈值的竖线与同一概率的横线，对照 CDF 与分位数的排序。
\subsubsection{分位数与非负增量表示}
\par\medskip\noindent\begin{minipage}{\linewidth}
\textbf{图解：累积分布函数越高，分位数反而越小。}
记 $Q_i(p)=F_i^{-1}(p)=\inf\{x:F_i(x)\ge p\}$，$0<p<1$。
这里 $F_i^{-1}$ 表示\textbf{广义反函数（分位数函数）}，不是 $1/F_i$；连续且严格递增时就是通常的反函数。

\begin{center}
\begin{minipage}{\linewidth}
\centering
\begin{minipage}[t]{0.49\linewidth}
\centering
\textbf{同一阈值：竖着比较概率}\\[3pt]
\resizebox{\linewidth}{!}{%
\WJInputFigure{09_累积分布函数比较}}
\par\smallskip
$F_B(x_0)>F_A(x_0)$
\end{minipage}\hfill
\begin{minipage}[t]{0.49\linewidth}
\centering
\textbf{同一概率：横着比较收益}\\[3pt]
\resizebox{\linewidth}{!}{%
\WJInputFigure{09_分位数函数比较}}
\par\smallskip
$F_B^{-1}(p)<F_A^{-1}(p)$
\end{minipage}
\end{minipage}
\end{center}
\end{minipage}\par\medskip

左图固定 $x_0$：沿竖直辅助线，$F_B$ 在上，表示 $B$ 不超过该阈值的概率更大。
右图固定 $p$：沿水平辅助线，先遇到 $F_B$、后遇到 $F_A$，再向下投影到横轴，故达到同一累计概率时，$B$ 的收益分位数更小。
图中用严格不等式展示两条曲线分开的情形；一般的一阶占优判据为
\[
 F_B(x)\ge F_A(x)\quad(\forall x)
 \quad\Longleftrightarrow\quad
 F_B^{-1}(p)\le F_A^{-1}(p)\quad(\forall p\in(0,1)).
\]

\begin{solution}[label-content={证明},blank-type=none]
\textbf{效用判据与分布判据。}先设收益取值有界，取 $L$ 严格低于两分布的支撑下端，$H$ 不低于两分布的支撑上端。记 $\Delta F=F_A-F_B$，对光滑 $u$，先用分布函数表示期望，再分部积分：
\[
\begin{aligned}
 \mathrm E[u(\tilde x_A)]-\mathrm E[u(\tilde x_B)]
 &=\int_L^H u(x)\,d\Delta F(x)\\
 &=\underbrace{[u(x)\Delta F(x)]_L^H}_{0}
   -\int_L^H u'(x)\Delta F(x)\,dx.
\end{aligned}
\]
边界项为零：两分布在支撑下端之前均为零，在上端均为一。

若 $\Delta F\le0$，右端对一切 $u'\ge0$ 非负，得到 FSD。反过来，若所有递增效用都偏好 $A$，取阈值效用 $u_t(x)=\mathbf1_{\{x>t\}}$（或其光滑递增近似），则
\[
 1-F_A(t)=\mathrm E[u_t(\tilde x_A)]
 \ge\mathrm E[u_t(\tilde x_B)]=1-F_B(t),
\]

所以 $F_A(t)\le F_B(t)$。

\textbf{非负增量表示。}由广义反函数的定义和 $F_i$ 的右连续性，$Q_i(p)\le x$ 当且仅当 $p\le F_i(x)$。令 $U$ 在 $(0,1)$ 上均匀分布，则
\[
 \Pr(Q_i(U)\le x)=\Pr(U\le F_i(x))=F_i(x).
\]

因此 $Q_i(U)$ 与原收益率 $\tilde x_i$ 同分布，分布函数有跳跃时这个结论也成立。用同一个 $U$ 为两只证券配对，由 $F_A\le F_B$ 得
\[
 \tilde x_A'=Q_A(U),\qquad \tilde x_B'=Q_B(U),\qquad
 \tilde d=Q_A(U)-Q_B(U)\ge0,
\]
这就得到第 3 条。反过来，由效用的单调性，$u(\tilde x_B'+\tilde d)\ge u(\tilde x_B')$，取期望即得一阶随机占优。
\end{solution}

一阶随机占优判据也适用于离散分布。

\textbf{经济含义：}取 $u(x)=x$ 得 $\mathrm E[\tilde x_A]\ge\mathrm E[\tilde x_B]$，但均值更高本身不能推出 FSD；还须在\textbf{每一个阈值}上，$A$ 落到该阈值以下的概率都不大于 $B$。若分布函数交叉，一阶判据不能给出排序。

\subsubsection{与无风险证券组合的福利比较}
\begin{tcolorbox}[knowledge,breakable=false]
\WJNineTheorem{thm:wj-nine-2}{} 若 $A\succeq_{\mathrm{FSD}}B$ 且 $\mathrm E[\tilde r_B]\ge r_F$，则对 $u'>0,u''<0$ 的参与者，
\[
 \max_a\mathrm E[u(b+a(\tilde r_A-r_F))]
 \ge\max_a\mathrm E[u(b+a(\tilde r_B-r_F))],
 \qquad b=w(1+r_F).
\]
两侧分别在一只风险证券与无风险证券之间配置，沿用\WJRef{sec:wj-eight-single}{第~}[~节（单只风险证券）]的自由交易和内点条件。把 $B$ 的最优投资金额用在 $A$ 上时，所得组合也须可行。
\end{tcolorbox}
\begin{solution}[label-content={证明},blank-type=none]
令 $a_B^*$ 为只交易 $B$ 与无风险证券时的最优风险投资额。由\WJRef{thm:wj-eight-2}{定理~}[（最优风险投资额与风险溢价同号）]，$\mathrm E[\tilde r_B]\ge r_F$ 保证 $a_B^*\ge0$\footnote{第八讲允许借贷和卖空，要求零头寸可行，且最优解是可双向微调的内点。对非退化收益，严格凹效用使目标的导数递减；在零头寸处，导数为 $u'(b)(\mathrm E[\tilde r_B]-r_F)$。因此风险溢价非负时，导数的零点不会在零头寸左侧。}，因此 $t\mapsto u(b+a_B^*(t-r_F))$ 非递减。先用同一金额把 $B$ 换成 $A$，由一阶随机占优，期望效用至少保持原水平。再对 $A$ 的投资额重新优化，得到
\[
\begin{aligned}
 \max_a\mathrm E[u(b+a(\tilde r_A-r_F))]
 &\ge\underbrace{\mathrm E[u(b+a_B^*(\tilde r_A-r_F))]}_{\text{沿用 }B\text{ 的最优投资额买 }A}\\
 &\ge\mathrm E[u(b+a_B^*(\tilde r_B-r_F))].
\end{aligned}
\]

最后一项就是只交易 $B$ 与无风险证券时的最优效用。
\end{solution}

定理比较两种投资机会各自能达到的最高效用。最优风险投资额还取决于收益分布与边际效用的相互作用。与第三只风险证券搭配时，组合风险取决于各证券在同一状态下的收益，见\WJRef{ex:wj-nine-2}{例~} 后的比较。

\subsection{二阶随机占优}
\begin{tcolorbox}[knowledge,breakable=false]
\WJNineDefinition{def:wj-nine-2}{（本教材口径）} 若对每个\textbf{\WJLink{def:wj-seven-1}{凹效用函数}}（$u''\le0$）都有
\[
 \mathrm E[u(\tilde x_A)]\ge\mathrm E[u(\tilde x_B)],
\]
则称 $A$ 对 $B$ \textbf{二阶随机占优}（SSD），记为 $A\succeq_{\mathrm{SSD}}B$。
\end{tcolorbox}
\textbf{凹性允许边际效用取正值或负值。}分别代入 $u(x)=x$ 和 $u(x)=-x$，两者都是凹函数，可得
\[
 \mathrm E[\tilde x_A]=\mathrm E[\tilde x_B].
\]

所以本教材的 SSD 是\textbf{等均值下的风险比较}，也可理解为凹序。另一些教材把 SSD 定义为所有\textbf{递增且凹}效用下的排序，那一定义容许均值不同；本章沿用王江教材的等均值口径。

\begin{tcolorbox}[knowledge,breakable=false]
\WJNineTheorem{thm:wj-nine-3}{} 下列叙述等价：
\begin{enumerate}
 \item $A\succeq_{\mathrm{SSD}}B$；
 \item $\mathrm E[\tilde x_A]=\mathrm E[\tilde x_B]$，且
 \[
 S(y):=\int_{-\infty}^{y}[F_A(x)-F_B(x)]\,dx\le0,\qquad\forall y;
 \]
 \item 存在同分布构造，使 $\tilde x_B\sim_d\tilde x_A+\tilde e$，且
 $\mathrm E[\tilde e\mid\tilde x_A]=0$。
\end{enumerate}
\end{tcolorbox}
第 3 条表示 $B$ 的分布可以由 $A$ 加上一层\textbf{条件均值为零的风险}得到，称为均值保持扩散。给定 $A$ 的每个收益水平，这层风险都是\WJRef{def:wj-seven-2}{定义~}[（公平赌博）]所说的公平赌博。这里的公平性要在每个收益水平分别检验。残差可以依赖 $A$ 的收益水平，例如在不同收益水平上有不同的波动幅度。

对任意收益目标 $y$，$(y-\tilde x)_+$ 表示收益缺口，达到目标时缺口为零。给定一次收益实现，介于该收益与目标 $y$ 之间的阈值区间长度就是缺口。先对阈值积分，再对收益取期望，可得
\[
 \mathrm E[(y-\tilde x)_+]
 =\mathrm E\!\left[\int_{-\infty}^y \mathbf1_{\{\tilde x\le t\}}\,dt\right]
 =\int_{-\infty}^y \Pr(\tilde x\le t)\,dt
 =\int_{-\infty}^y F(t)\,dt.
\]

因此 $S(y)$ 是 $A$ 的平均缺口减去 $B$ 的平均缺口。本教材的 SSD 要求均值相同，且 $A$ 在每个目标下的平均缺口都不大于 $B$。\WJRef{fig:wj-mean-preserving-spread}{图~} 展示了 CDF 交叉而 $S(y)$ 始终非正的情形。

\WJInputFigure{09_均值保持扩散}

\begin{solution}[label-content={证明},blank-type=none]
\textbf{1 与 2 等价。}仍先考虑有界支撑 $[L,H]$。等均值给出
\[
 S(L)=0,\qquad S(H)=\mathrm E[\tilde x_B]-\mathrm E[\tilde x_A]=0.
\]

连续分部积分两次，得到
\[
\begin{aligned}
 \mathrm E[u(\tilde x_A)]-\mathrm E[u(\tilde x_B)]
 &=-\int_L^H u'(x)[F_A(x)-F_B(x)]\,dx\\
 &=\underbrace{-[u'(x)S(x)]_L^H}_{0}+\int_L^H u''(x)S(x)\,dx.
\end{aligned}
\]

$S\le0$、$u''\le0$ 时右端非负，故 2 推出 1。反向，除利用线性效用得到等均值外，对每个 $y$ 取凹函数 $u_y(x)=\min\{x,y\}$。由于
\[
 \mathrm E[\min\{\tilde x,y\}]
 =y-\mathrm E[(y-\tilde x)_+],\qquad
 \mathrm E[(y-\tilde x)_+]=\int_{-\infty}^y F(x)\,dx,
\]

所有凹效用均偏好 $A$ 就要求 $-S(y)\ge0$，得到 2。非光滑函数可用凹的光滑函数近似。

\textbf{3 推出 1。}先固定 $\tilde x_A$ 的取值，只对附加风险 $\tilde e$ 取条件平均。由于 $u$ 是凹函数，先对附加风险取条件均值再计算效用，所得效用至少与先计算效用再取条件平均一样高。这就是这里用到的条件 Jensen 不等式。再对 $\tilde x_A$ 取期望，并用 $\mathrm E[\tilde e\mid\tilde x_A]=0$，得到
\[
\begin{aligned}
 \mathrm E[u(\tilde x_B)]
 &=\mathrm E\{\mathrm E[u(\tilde x_A+\tilde e)\mid\tilde x_A]\}\\
 &\le\mathrm E\{u(\tilde x_A+\mathrm E[\tilde e\mid\tilde x_A])\}
 =\mathrm E[u(\tilde x_A)].
\end{aligned}
\]

\textbf{1 推出 3。}这一步用到均值保持扩散表示定理（Rothschild--Stiglitz，1970）。该定理保证可以把两个收益率的分布放到同一个联合分布中，使 $\tilde x_A+\tilde\varepsilon$ 与 $\tilde x_B$ 同分布，且 $\mathrm E[\tilde\varepsilon\mid\tilde x_A]=0$。
\end{solution}

给定扩散前的收益率，扩散后的条件均值仍等于该收益率。

\textbf{方差的作用：}若二阶矩存在，取 $u(x)=-x^2$，SSD 推出
$\operatorname{Var}(\tilde x_A)\le\operatorname{Var}(\tilde x_B)$。反之，等均值加方差更小并不足以保证 SSD，因为凹效用还会反映尾部等分布特征。

\subsection{随机占优与逐状态占优}
\begin{tcolorbox}[knowledge,breakable=false]
\WJNineDefinition{def:wj-nine-3}{} 若在实际同一组状态下，$x_A(\omega)\ge x_B(\omega)$ 对每个 $\omega$ 都成立，且至少一个可能状态严格大于，则称 $A$ \textbf{占优}于 $B$。
\end{tcolorbox}
逐状态占优比较同一状态下的表现，不依赖各状态概率的具体大小；随机占优比较\textbf{概率分布}，依赖参与者采用的概率。逐状态占优推出 FSD，FSD 不推出实际逐状态占优。

若每投入 1 元，$A$ 在所有状态的总支付不低于 $B$，某状态严格更高，则买入 1 元 $A$、卖空 1 元 $B$ 的初始净成本为零，未来净支付处处非负且在正概率状态为正，构成套利；这里假设允许卖空、按相同单位投资额交易且没有交易摩擦。随机占优本身不能保证这个逐状态支付差，因此不必产生套利，随机被占优的证券仍可能存在于均衡中。

\begin{tcolorbox}[example,breakable=false]
\begin{question}
\WJNineExample{ex:wj-nine-1}{（逐状态占优）} 两种状态下的总收益率如下：
\begin{center}
\begin{tabular}{c|cc}
 &好状态&坏状态\\\hline
 $A$&1.2&1.0\\
 $B$&1.1&1.0
\end{tabular}
\end{center}
判断 $A$ 是否逐状态占优及一阶随机占优于 $B$。
\end{question}
\begin{solution}[blank-type=manual,blank-vsep=3\baselineskip]
$A$ 在好状态严格更好，在坏状态相同，所以 $A$ 占优于 $B$，并一阶随机占优于 $B$。
\end{solution}
\end{tcolorbox}

\begin{tcolorbox}[example,breakable=false]
\begin{question}
\WJNineExample{ex:wj-nine-2}{（FSD 不等于逐状态占优）} 好、坏状态等概率，总收益率为
\begin{center}
\begin{tabular}{c|cc}
 &好状态&坏状态\\\hline
 $A$&1.2&1.0\\
 $B$&1.0&1.1
\end{tabular}
\end{center}
判断 $A$ 是否逐状态占优及一阶随机占优于 $B$，并说明是否构成套利。
\end{question}
\begin{solution}[blank-type=manual,blank-vsep=3\baselineskip]
$A$ 在坏状态低于 $B$，故不逐状态占优。其分布却满足
\[
\begin{array}{c|cccc}
 t&t<1.0&1.0\le t<1.1&1.1\le t<1.2&t\ge1.2\\\hline
 F_A(t)&0&1/2&1/2&1\\
 F_B(t)&0&1/2&1&1
\end{array}
\]

处处 $F_A\le F_B$，所以 $A\succeq_{\mathrm{FSD}}B$。但买 $A$、卖 $B$ 的状态收益差为 $[0.2;-0.1]$，有亏损状态，不能据此认定套利。
\end{solution}
\end{tcolorbox}

% 教学补充：用教材例 9.2 的数据具体说明印刷第 127--128 页关于组合的讨论。
沿用\WJRef{ex:wj-nine-2}{例~}，加入总收益率为 $(0.8,1.3)$ 的证券 $C$，两项依次对应好、坏状态。分别将 $A$、$B$ 与 $C$ 各按一半权重搭配，得到 $(1.0,1.15)$ 和 $(0.9,1.20)$。以 $0.95$ 为阈值，第一个组合落在阈值以下的概率为零，第二个为 $1/2$。以 $1.17$ 为阈值，两者的概率分别为 $1$ 和 $1/2$。两条 CDF 因而交叉，原有的一阶占优关系消失。\WJRef{exer:wj-nine-2}{习题~} 进一步说明，即使保持各资产的边际分布不变，改变联合分布也可能改变最优持仓。

\section{组合分离}\label{sec:wj-nine-2}
随机占优可以排序给定的组合，前面的例子也表明，引入其他风险证券会改变这种排序。现在换一个问题，考察收益分布或偏好具有什么共同结构时，参与者可以用同一组基金构造各自的最优组合。

\subsection{共同基金分离}
令 $w>0$ 为总投资额，$a_n$ 为第 $n$ 只证券的投资金额。组合用\textbf{权重列向量} $z=[z_1,\ldots,z_N]^{\mathsf T}$ 表示，其中 $z_n=a_n/w$。按同一证券顺序记总收益率列向量 $\tilde x=[\tilde x_1,\ldots,\tilde x_N]^{\mathsf T}$，并令 $\iota$ 为 $N$ 维全 $1$ 列向量。组合权重之和为 $1$，组合总收益率与期末财富分别为
\[
 z^{\mathsf T}\iota=1,\qquad
 \tilde x_z=z^{\mathsf T}\tilde x=\sum_{n=1}^N z_n\tilde x_n,\qquad
 \tilde w=w\tilde x_z.
\]
给定 $w$ 后令 $v(t)=u(wt)$，问题为
\[
 \max_{z:z^{\mathsf T}\iota=1}\mathrm E[v(\tilde x_z)].
\]
$v$ 既受原效用 $u$ 影响，也受财富 $w$ 影响。因此不同财富者即使有相同的财富效用，也未必选择相同权重。以下采用共同概率信念，允许自由卖空与借贷，并讨论最优解存在的情形。用纯风险基金表示共同风险方向时，该方向的净投资额须非零。

\begin{tcolorbox}[knowledge,breakable=false]
\WJNineDefinition{def:wj-nine-4}{} 若存在 $F$ 只共同基金 $z_{f_1},\ldots,z_{f_F}$，使每个参与者都能选择一个最优组合表示为
\begin{equation}
 z_k=\sum_{i=1}^F\lambda_{ki}z_{f_i},\qquad
 \sum_{i=1}^F\lambda_{ki}=1.\label{eq:wj-9-1}
\end{equation}
则称 $F$--\textbf{基金分离}成立。
\end{tcolorbox}
最优解唯一时，这些最优权重位于至多 $F-1$ 维的\textbf{仿射}空间中；它由 $F$ 只基金张成，且所有权重之和仍为 1。$\lambda_{ki}$ 可以为负，故不是只能取非负系数的凸组合。基金由原证券构造，分离基金的选择通常不唯一。若最优解不唯一，分离要求每人有一个最优选择可由共同基金实现；覆盖全部最优组合则需另行考察其集合维数。

$z_{f_i}$ 的各分量是基金 $i$ 投在原证券上的权重，$\lambda_{ki}$ 是参与者 $k$ 投在该基金上的财富比例。因此，参与者 $k$ 对基金 $i$ 的投资金额为 $w_k\lambda_{ki}$。按基金把各人的金额加总后，整个市场对原证券的投资需求也由同一组基金权重向量生成。这就把个体组合的共同结构与市场总需求联系起来。

\begin{tcolorbox}[knowledge,breakable=false]
\WJNineDefinition{def:wj-nine-5}{} 若上述 $F$ 只基金中有一只为无风险证券，其余 $F-1$ 只只由风险证券组成，则称 $F$--\textbf{基金货币分离}成立。
\end{tcolorbox}
两基金货币分离意味着每个人都在“同一只风险基金”和“无风险证券”之间分配财富。不同人的风险承受程度可以改变两者的比例，但不改变风险基金内部的相对权重。$N$ 基金分离总可用原来 $N$ 只证券实现，有内容的结论是 $F<N$。

\subsection{收益率分布与共同基金分离}
\begin{tcolorbox}[knowledge,breakable=false]
% 按课堂投影片保留定理措辞、三条条件及权重符号 w_n。
\WJNineTheorem{thm:wj-nine-4}{（Ross）} 假设存在总收益率为 $x_F$ 的无风险证券。考虑以下三条因子条件。对每只风险证券 $n=1,\ldots,N-1$，存在载荷 $\beta_n$、基金权重 $w_n$ 和残差 $\tilde e_n$，以及共同因子 $\tilde y$，使得
\begin{enumerate}[ref=\arabic*]
 \item $\displaystyle \tilde x_n=x_F+\beta_n\tilde y+\tilde e_n$；
 \item\label{item:wj-nine-ross-zero-conditional-mean} $\displaystyle \mathrm E[\tilde e_n\mid\tilde y]=0$；
 \item\label{item:wj-nine-ross-zero-fund-residual} $\displaystyle \sum_{n=1}^{N-1}w_n\tilde e_n=0$，其中
       $\displaystyle \sum_{n=1}^{N-1}w_n=1$。
\end{enumerate}
其中风险基金因子载荷 $\beta_w=\sum_{n=1}^{N-1}w_n\beta_n\ne0$。

\textbf{充分性。}满足上述条件时，每个有最优解的递增凹效用投资者，都能用权重为 $(w_1,\ldots,w_{N-1})$ 的风险基金与无风险证券构成一个最优组合。

\textbf{必要性（有限状态）。}若存在对上述投资者共同适用的风险基金，各状态概率为正、该基金的期望超额收益非零，且有一种光滑的收益率效用 $v_0$ 具有非零的内点最优风险头寸，并在相关区间上满足 $v_0'>0$、$v_0''<0$，则可以反推出这三条因子条件。
\end{tcolorbox}

\textbf{符号说明：}这里 $x_F=1+r_F$ 是无风险总收益率，$w_n$ 是风险基金在第 $n$ 只风险证券上的权重。用 $w=[w_1,\ldots,w_{N-1},0]^{\mathsf T}$ 表示风险基金的权重向量，用 $F=[0,\ldots,0,1]^{\mathsf T}$ 表示无风险基金；本段的 $w$ 是向量，初始财富另记为 $W$。无风险证券作为分离基金之一，故这里也属于\WJRef{def:wj-nine-5}{定义~}[（基金货币分离）] 的两基金货币分离。

\textbf{条件含义：}共同因子产生\textbf{共同风险}，残差产生\textbf{特殊风险}。条件均值为零约束残差在每个因子取值上的平均值；第 3 条要求可交易风险基金\textbf{完全消除}残差。基金的非零因子载荷使任意组合的共同风险都可由该基金匹配。

例如，两只证券相对无风险证券的超额收益分别为 $\tilde y+\tilde e$ 和 $\tilde y-\tilde e$，且 $\mathrm E[\tilde e\mid\tilde y]=0$。取 $w=[1/2,1/2]$、$\beta=[1,1]$，则 $\beta_w=1$，各投一半所得基金的超额收益恰为 $\tilde y$。任意权重 $[\theta,1-\theta]$ 的超额收益为
\[
 \underbrace{\tilde y}_{\text{共同因子}}
 +\underbrace{(2\theta-1)\tilde e}_{\text{剩余风险}}.
\]
给定 $\tilde y$ 后剩余风险的条件均值为零，因此凹效用下去掉它不会降低期望效用。

% 课堂照片补充：投影片“B 收益率分布与共同基金分离”，定理~\ref{thm:wj-nine-4}；
% 板书由 E[Y|X]=0 推出 E[Y]=0，并展开 Cov(X,Y)。以下补齐中间步骤。
\subsubsection{条件均值、协方差与因子载荷}
设随机变量 $X,Y$ 的二阶矩有限，且 $\mathrm E[Y\mid X]=0$（几乎处处）。先用迭代期望法则求 $Y$ 的无条件均值：
\[
\begin{aligned}
 \mathrm E[Y]
 &=\mathrm E\!\left[\mathrm E[Y\mid X]\right]\\
 &=\mathrm E[0]=0.
\end{aligned}
\]

协方差可展开为：
\[
\begin{aligned}
 \operatorname{Cov}(X,Y)
 &=\mathrm E\!\left[(X-\mathrm E[X])(Y-\mathrm E[Y])\right]\\
 &=\mathrm E[XY]-\mathrm E[X]\mathrm E[Y]
   -\mathrm E[Y]\mathrm E[X]+\mathrm E[X]\mathrm E[Y]\\
 &=\mathrm E[XY]-\mathrm E[X]\mathrm E[Y].
\end{aligned}
\]
其中乘积项也可用迭代期望计算。由于给定 $X$ 后，$X$ 本身已知，可以提出条件期望：
\[
\begin{aligned}
 \mathrm E[XY]
 &=\mathrm E\!\left[\mathrm E[XY\mid X]\right]\\
 &=\mathrm E\!\left[X\,\mathrm E[Y\mid X]\right]\\
 &=\mathrm E[X\cdot0]=0.
\end{aligned}
\]

将两项代回，便得到
\[
 \boxed{\operatorname{Cov}(X,Y)=0-\mathrm E[X]\cdot0=0.}
\]

若 $X,Y$ 的方差都为正，相应的相关系数也为零；零方差情形的协方差仍为零。

\textbf{代回\WJRef{thm:wj-nine-4}{定理~}[（Ross 两基金分离定理）]：}令 $X=\tilde y$、$Y=\tilde e_n$，便有
\[
 \mathrm E[\tilde e_n]=0,\qquad
 \mathrm E[\tilde y\tilde e_n]=0,\qquad
 \operatorname{Cov}(\tilde y,\tilde e_n)=0.
\]

因此对收益结构 $\tilde x_n=x_F+\beta_n\tilde y+\tilde e_n$，
\[
\begin{aligned}
 \mathrm E[\tilde x_n]&=x_F+\beta_n\mathrm E[\tilde y],\\
 \operatorname{Cov}(\tilde x_n,\tilde y)
 &=\beta_n\operatorname{Var}(\tilde y)
   +\operatorname{Cov}(\tilde e_n,\tilde y)
 =\beta_n\operatorname{Var}(\tilde y).
\end{aligned}
\]

当 $\operatorname{Var}(\tilde y)>0$ 时，
\[
 \beta_n=\frac{\operatorname{Cov}(\tilde x_n,\tilde y)}
                   {\operatorname{Var}(\tilde y)}.
\]

更一般地，只要乘积可积，对任意 $g(\tilde y)$ 都有
\[
 \mathrm E[g(\tilde y)\tilde e_n]
 =\mathrm E\!\left[g(\tilde y)\mathrm E[\tilde e_n\mid\tilde y]\right]=0.
\]

这说明条件均值为零还排除了残差与因子的非线性函数之间的这类联系。它不要求残差与因子独立，也不要求不同证券的残差彼此不相关。

\textbf{逆命题不成立：}令 $X$ 以相同概率取 $-1,0,1$，$Y=X^2-2/3$。则
\[
 \mathrm E[X]=\mathrm E[Y]=0,\qquad
 \mathrm E[XY]=\mathrm E[X^3]-\tfrac23\mathrm E[X]=0,
\]

相应的 $Y$ 依次为 $1/3,-2/3,1/3$，所以 $\operatorname{Cov}(X,Y)=0$；但 $Y$ 是 $X$ 的函数，故
\[
 \mathrm E[Y\mid X]=X^2-\tfrac23\ne0.
\]

因此，定理中的条件均值限制不能仅用“不相关”替代。

% 课堂照片补充：投影片第 27 页“充分性”及板书的向量展开。
\subsubsection{Ross 定理的充分性}
风险基金能匹配原组合对共同因子的暴露，同时去掉条件均值为零的剩余风险。

\begin{solution}[label-content={证明},blank-type=none]
\textbf{第一步：将各证券的收益写成向量形式。}
令 $\iota=[1,\ldots,1]^{\mathsf T}$，并设
\[
\begin{gathered}
 \tilde x=[\tilde x_1,\ldots,\tilde x_{N-1},x_F]^{\mathsf T},\\
 \beta=[\beta_1,\ldots,\beta_{N-1},0]^{\mathsf T},\qquad
 \tilde e=[\tilde e_1,\ldots,\tilde e_{N-1},0]^{\mathsf T}.
\end{gathered}
\]

最后一只证券是无风险证券，故其因子载荷和残差都为零。逐只证券的收益等式合为
\[
 \boxed{\tilde x=x_F\iota+\beta\tilde y+\tilde e.}
\]

\textbf{第二步：构造只有共同风险的基金 $w$。}
由\WJLink{item:wj-nine-ross-zero-fund-residual}{定理~\ref*{thm:wj-nine-4}~第~\ref*{item:wj-nine-ross-zero-fund-residual}~条（风险基金权重之和为 $1$，且组合残差为零）}，$w^{\mathsf T}\iota=1$、$w^{\mathsf T}\tilde e=0$。定义 $\beta_w=w^{\mathsf T}\beta$，于是
\[
\begin{aligned}
 \tilde x_w:=w^{\mathsf T}\tilde x
 &=x_Fw^{\mathsf T}\iota+w^{\mathsf T}\beta\tilde y
   +w^{\mathsf T}\tilde e\\
 &=x_F+\beta_w\tilde y.
\end{aligned}
\]
基金 $w$ 的收益只包含共同风险。若初始财富为 $W>0$，全部投资于该基金后的期末财富为
\[
 \tilde W=Ww^{\mathsf T}\tilde x=W\tilde x_w,
 \qquad \frac{\tilde W}{W}=\tilde x_w.
\]

因此，权重向量 $w$、基金总收益率 $\tilde x_w$ 与期末财富 $\tilde W$ 是三个不同的量。

\textbf{第三步：将任意组合分解为两基金组合与剩余头寸。}
任取满足 $z^{\mathsf T}\iota=1$ 的组合 $z$，记 $\beta_z=z^{\mathsf T}\beta$。在 $\beta_w\ne0$ 时，令
\[
 \lambda=\frac{\beta_z}{\beta_w}
 =\frac{z^{\mathsf T}\beta}{w^{\mathsf T}\beta},\qquad
 z_0=\lambda w+(1-\lambda)F,\qquad \eta=z-z_0.
\]
便得到三部分分解
\[
 z=\underbrace{\lambda w+(1-\lambda)F}_{\text{匹配共同因子敞口的两基金组合}}
  +\underbrace{\eta}_{\text{零成本剩余头寸}}.
\]
这里 $\lambda$ 可为任意实数。由于 $w^{\mathsf T}\iota=F^{\mathsf T}\iota=1$，
\[
\begin{aligned}
 1=z^{\mathsf T}\iota
 &=\lambda w^{\mathsf T}\iota+(1-\lambda)F^{\mathsf T}\iota
   +\eta^{\mathsf T}\iota\\
 &=1+\eta^{\mathsf T}\iota,
 \qquad\Longrightarrow\qquad \eta^{\mathsf T}\iota=0.
\end{aligned}
\]

又因 $F^{\mathsf T}\beta=0$，所以上述分解两边乘以 $\beta$ 得
\[
\begin{aligned}
 z^{\mathsf T}\beta
 &=\lambda w^{\mathsf T}\beta+(1-\lambda)F^{\mathsf T}\beta
   +\eta^{\mathsf T}\beta\\
 &=\frac{\beta_z}{\beta_w}\beta_w+0+\eta^{\mathsf T}\beta\\
 &=z^{\mathsf T}\beta+\eta^{\mathsf T}\beta,
 \qquad\Longrightarrow\qquad \eta^{\mathsf T}\beta=0.
\end{aligned}
\]

所以 $\eta$ 既是\textbf{零成本头寸}，又没有共同因子敞口。

\textbf{第四步：展开组合收益，识别剩余风险。}
无风险基金有 $F^{\mathsf T}\tilde x=x_F$；而剩余头寸满足
\[
\begin{aligned}
 \eta^{\mathsf T}\tilde x
 &=x_F\eta^{\mathsf T}\iota+\eta^{\mathsf T}\beta\tilde y
   +\eta^{\mathsf T}\tilde e\\
 &=\eta^{\mathsf T}\tilde e.
\end{aligned}
\]

因此，组合收益可写为
\[
\begin{aligned}
 \tilde x_z=z^{\mathsf T}\tilde x
 &=\lambda w^{\mathsf T}\tilde x+(1-\lambda)F^{\mathsf T}\tilde x
   +\eta^{\mathsf T}\tilde x\\
 &=\lambda(x_F+\beta_w\tilde y)+(1-\lambda)x_F
   +\eta^{\mathsf T}\tilde e\\
 &=\underbrace{x_F+\beta_z\tilde y}_{\tilde x_{z_0}}
   +\eta^{\mathsf T}\tilde e.
\end{aligned}
\]

又由\WJLink{item:wj-nine-ross-zero-conditional-mean}{定理~\ref*{thm:wj-nine-4}~第~\ref*{item:wj-nine-ross-zero-conditional-mean}~条（给定共同因子时，各证券残差的条件均值为零）}，
\[
 \mathrm E[\eta^{\mathsf T}\tilde e\mid\tilde y]
 =\sum_{n=1}^{N-1}\eta_n\mathrm E[\tilde e_n\mid\tilde y]=0,
 \qquad
 \mathrm E[\tilde x_z\mid\tilde y]=\tilde x_{z_0}.
\]

\textbf{第五步：利用凹效用移除剩余风险。}
在自由交易下，可用风险基金和无风险证券组成 $z_0=\lambda w+(1-\lambda)F$。其收益是原组合收益给定 $\tilde y$ 的条件均值，因此仍落在 $v$ 的区间定义域内。给定 $\tilde y$ 后应用条件 Jensen 不等式：
\[
\begin{aligned}
 \mathrm E[v(\tilde x_z)\mid\tilde y]
 &\le v\!\left(\mathrm E[\tilde x_z\mid\tilde y]\right)
 =v(\tilde x_{z_0}),\\
 \mathrm E[v(\tilde x_z)]
 &=\mathrm E\!\left[\mathrm E[v(\tilde x_z)\mid\tilde y]\right]
 \le\mathrm E[v(\tilde x_{z_0})].
\end{aligned}
\]

由于 $\tilde x_{z_0}$ 是 $\tilde y$ 的函数，迭代期望给出
\[
 \mathrm E[\eta^{\mathsf T}\tilde e\mid\tilde x_{z_0}]
 =\mathrm E\!\left[\mathrm E[\eta^{\mathsf T}\tilde e\mid\tilde y]
                   \mid\tilde x_{z_0}\right]=0.
\]

所以原组合收益 $\tilde x_z$ 正是 $\tilde x_{z_0}$ 的均值保持扩散，与\WJRef{thm:wj-nine-3}{定理~}[（二阶随机占优的等价判据）] 的判据一致。把原组合换成保留同样因子敞口的 $z_0$ 后，期望效用至少保持原水平。因此可在 $w$、$F$ 中找到最优组合。
\end{solution}

若效用严格凹且剩余风险非退化，移除该风险会严格提高效用。弱凹或冗余时最优头寸未必唯一。

\subsubsection{Ross 定理的必要性（有限状态）}
% 教材印刷页 135 的勘误要求保留 x_F+lambda*y；以下以固定最优点的边际效用扰动补足量词。
在有限状态模型中，设各状态概率为正，风险基金超额收益 $\tilde y=\tilde x_w-x_F$ 满足 $\mathrm E[\tilde y]\ne0$。假设存在一个光滑收益率效用 $v_0$，在相关财富区间上满足 $v_0'>0$、$v_0''<0$，其两基金最优组合为
\[
 \tilde x_0=x_F+\lambda_0\tilde y,\qquad \lambda_0\ne0,
\]
且该解位于可行域内部。这些条件给出下面必要性证明的非退化情形。

\begin{solution}[label-content={证明},blank-type=none]
\textbf{第一步：构造均值为零的残差。}令
\[
 \beta_n=\frac{\mathrm E[\tilde x_n]-x_F}{\mathrm E[\tilde y]},\qquad
 \tilde e_n=\tilde x_n-x_F-\beta_n\tilde y.
\]

于是 $\mathrm E[\tilde e_n]=0$。对基金各证券加权，利用 $\sum_n w_n=1$ 与 $\tilde x_w-x_F=\tilde y$，得到
\[
 \sum_n w_n\beta_n
 =\frac{\mathrm E[\tilde x_w]-x_F}{\mathrm E[\tilde y]}=1,
 \qquad
 \sum_n w_n\tilde e_n
 =\tilde x_w-x_F-\tilde y\sum_nw_n\beta_n=0.
\]

\textbf{第二步：写出最优点的一阶条件。}两基金分离使 $\tilde x_0$ 同时为全市场最优收益率。分别沿基金与单只证券的超额收益方向求导，有
\[
 \mathrm E[v_0'(\tilde x_0)\tilde y]=0,\qquad
 \mathrm E[v_0'(\tilde x_0)(\tilde x_n-x_F)]=0.
\]

第二式减去第一式的 $\beta_n$ 倍，得
\[
 \mathrm E[v_0'(\tilde x_0)\tilde e_n]=0.
\]

\textbf{第三步：保持同一最优点，改变边际效用。}任取定义在 $\tilde y$ 各取值上的函数 $h$，满足 $\mathrm E[h(\tilde y)\tilde y]=0$。由于 $\lambda_0\ne0$，不同的 $y$ 对应不同的 $x_F+\lambda_0 y$。在这些点周围取互不相交的小邻域。对每一点取一个光滑函数，使其在该点等于 $1$、在邻域外等于 $0$。用各点的 $h(y)$ 加权这些函数，将其和乘以 $\varepsilon$ 后加到 $v_0'$ 上，得到
\[
 v_\varepsilon'(x_F+\lambda_0 y)
 =v_0'(x_F+\lambda_0 y)+\varepsilon h(y),
\]
各邻域可取在 $v_0'>0$、$v_0''<0$ 的区间内。由连续性，在这些小闭邻域上取充分小的 $\lvert\varepsilon\rvert$，可使 $v_\varepsilon'>0$、$v_\varepsilon''<0$ 继续成立。扰动后仍有
\[
 \mathrm E[v_\varepsilon'(\tilde x_0)\tilde y]=0.
\]

严格凹性使基金上的最优支付仍为 $\tilde x_0$。由共同基金分离的假设，这也是扰动后的全市场最优支付。像第二步一样，对 $v_\varepsilon$ 用证券方向的一阶条件减去基金方向的一阶条件的 $\beta_n$ 倍。扰动前后分别有
\[
 \mathrm E[v_0'(\tilde x_0)\tilde e_n]=0,\qquad
 \mathrm E[v_\varepsilon'(\tilde x_0)\tilde e_n]=0.
\]
两式相减，代入 $v_\varepsilon'(\tilde x_0)-v_0'(\tilde x_0)=\varepsilon h(\tilde y)$，再除以非零的 $\varepsilon$，得到 $\mathrm E[h(\tilde y)\tilde e_n]=0$。

这里通过保持基金方向的一阶条件，让效用扰动前后的最优投资量都等于 $\lambda_0$。

\textbf{第四步：识别条件均值。}记 $m_n(y)=\mathrm E[\tilde e_n\mid\tilde y=y]$。上式说明，所有满足 $\mathrm E[h(\tilde y)\tilde y]=0$ 的函数都满足
\[
 \mathrm E[h(\tilde y)m_n(\tilde y)]=0.
\]

令 $c_n=\mathrm E[m_n(\tilde y)\tilde y]/\mathrm E[\tilde y^2]$，取投影残差 $h(\tilde y)=m_n(\tilde y)-c_n\tilde y$。按 $c_n$ 的定义，$\mathrm E[h(\tilde y)\tilde y]=0$。把这个 $h$ 代入上式，有
\[
 0=\mathrm E[h(\tilde y)m_n(\tilde y)]
 =\mathrm E[h(\tilde y)^2]+c_n\mathrm E[h(\tilde y)\tilde y]
 =\mathrm E[h(\tilde y)^2].
\]
故 $h=0$，即 $m_n(\tilde y)=c_n\tilde y$。再取期望：
\[
 0=\mathrm E[\tilde e_n]=c_n\mathrm E[\tilde y].
\]

由 $\mathrm E[\tilde y]\ne0$ 得 $c_n=0$，故 $\mathrm E[\tilde e_n\mid\tilde y]=0$，三条因子条件全部成立。
\end{solution}
% 一般分布的必要性涉及函数空间与扰动可容许性；本讲展开有限状态、非退化内点情形。

\subsection{偏好与共同基金分离}
本节改为限制参与者偏好，而不预先限制证券收益率分布。
\subsubsection{单基金分离与共同收益率偏好}
\begin{tcolorbox}[knowledge,breakable=false]
\WJNineTheorem{thm:wj-nine-5}{} 在\textbf{任意证券集合}下均成立单基金分离，当且仅当所有参与者定义在\textbf{收益率}上的效用函数在正仿射变换意义下相同，即
\[
 v_k(t)=A_k+B_kv(t),\qquad B_k>0.
\]
\end{tcolorbox}
\begin{solution}[label-content={证明},blank-type=none]
\textbf{充分方向。}正仿射变换不改变期望效用的排序，故所有人都可以选择同一个最优权重组合，只是投资金额不同。
\end{solution}

这里的收益率效用为 $v_k(t)=u_k(w_kt)$，因此财富差异也进入偏好比较。反向是对任意证券集合成立的偏好刻画，依赖可行彩票足够丰富及期望效用表示的唯一性；这里作为外部刻画引用，不能从某一个给定市场中的单基金分离推出。

例如，两人都采用对数财富效用时，对任意正总收益率 $t$，$v_k(t)=\log(w_kt)=\log w_k+\log t$。初始财富 $w_k$ 只改变效用的常数项，所以两人对各组合的排序相同，可以选择同一组最优权重。

\subsubsection{二次效用与两基金分离}
\begin{tcolorbox}[knowledge,breakable=false]
\WJNineTheorem{thm:wj-nine-6}{} 在\WJRef{sec:wj-twelve-2}{第~}[~节（均值--方差前沿组合）]的市场设定下，参与者共享收益均值、协方差矩阵和可行组合集合，并允许自由交易。若每人的最优组合存在，且效用在相关财富区间上为严格递增、严格凹的二次函数，则两基金分离成立。
\end{tcolorbox}
% 教材印刷第 136 页将证明指向第 12 章，第 184 页把二次效用作为均值--方差两基金分离的特例。
\begin{solution}[label-content={证明},blank-type=none]
由\WJRef{thm:wj-twelve-1}{定理~}[（二次效用的均值--方差表示）]，期望效用只依赖期末财富的均值和方差。同均值下，方差项的系数为负，因此各人严格偏好较低方差。由\WJRef{thm:wj-twelve-4}{定理~}[（最优组合位于均值--方差前沿）]，所有人的最优组合都在同一条前沿上。\WJRef{cor:wj-twelve-1}{推论~}[（两个前沿组合生成整条前沿）] 说明任选两个不同的前沿组合即可生成整条前沿，\WJRef{thm:wj-twelve-6}{定理~}[（均值--方差两基金分离）] 据此得到两基金分离。
\end{solution}

这项结论也适用于只含风险证券的市场。

\textbf{教材的反向表述。}教材以充要形式陈述定理。按本讲\WJEqRef{eq:wj-9-1}{基金分离的组合表示~}，基金权重允许为零，单基金分离也满足两基金表示。例如，共同采用对数效用的参与者，在共同可行权重集合内有最优解时，可以选择同一只基金，并令第二只基金的权重为零。这个例子表明，在本讲定义下，两基金分离也可以出现在非二次效用中。

\subsubsection{HARA 偏好与两基金货币分离}
\begin{tcolorbox}[knowledge,breakable=false]
\WJNineTheorem{thm:wj-nine-7}{（Cass--Stiglitz，充分方向）} 存在总收益率 $x_F>0$ 的无风险证券。设参与者有共同信念，允许自由交易，全投无风险证券时的财富位于各自效用的定义域内，且各人的最优解为内点。若每位参与者的边际效用在正比例变换后均可写成
\[
 u_k'(t)=\left(d_k+\frac{t}{\gamma}\right)^{-\gamma},
 \qquad d_k+\frac{t}{\gamma}>0.
\]
其中 $\gamma\ge-1$、$\gamma\ne0$ 对所有参与者相同，则最优风险头寸可取为同一方向的倍数。当该共同方向的净投资额非零时，可将其归一化为一只纯风险基金，与无风险证券共同实现两基金货币分离。
\end{tcolorbox}
这是具有共同形状参数的 \WJLink{sec:wj-seven-hara}{HARA 偏好}。教材在定理中取 $d_k>0$，在充分性证明中也讨论 $d_k=0$。取 $\gamma>0$、$d_k=0$ 得到 CRRA 特例。负指数效用是风险容忍度恒定的极限情形。HARA 的绝对风险容忍度为
\[
 T_k(t)=-\frac{u_k'(t)}{u_k''(t)}=d_k+\frac{t}{\gamma}.
\]

关键是所有人的风险容忍度对财富具有\textbf{相同斜率}，截距可以不同。对数效用、幂效用、二次效用及负指数效用的相应形式或极限均与这一族有关；正的仿射变换不影响组合选择。

\begin{solution}[label-content={证明},blank-type=none]
假设共同信念、自由交易、全投无风险证券时的财富位于效用定义域内，且存在内点最优解。令 $\tilde x_R=[\tilde x_1,\ldots,\tilde x_{N-1}]^{\mathsf T}$ 为风险证券的总收益率列向量，此处 $\iota$ 是同维的全 $1$ 列向量。于是 $\tilde q=\tilde x_R-x_F\iota$ 的第 $n$ 项是第 $n$ 只风险证券的超额收益率。令 $a_k$ 为按同一顺序排列的风险证券投资\textbf{金额}列向量。期末财富为
\[
 \tilde w_k=w_kx_F+a_k^{\mathsf T}\tilde q.
\]

全投无风险证券时的期末财富为 $w_kx_F$。在该财富处的风险容忍度为 $C_k=T_k(w_kx_F)=d_k+w_kx_F/\gamma>0$。用它缩放风险投资金额，令
\[
 a_k=\gamma C_k h_k,
\]
则
\[
 d_k+\frac{\tilde w_k}{\gamma}
 =C_k(1+h_k^{\mathsf T}\tilde q).
\]

内点一阶条件 $\mathrm E[u_k'(\tilde w_k)\tilde q]=0$ 变成
\[
 C_k^{-\gamma}\mathrm E[(1+h_k^{\mathsf T}\tilde q)^{-\gamma}\tilde q]=0.
\]

消掉正的常数 $C_k^{-\gamma}$，方程及 HARA 自然定义域 $1+h_k^{\mathsf T}\tilde q>0$ 都与 $w_k,d_k$ 无关。由凹性，大家可以选择同一个解 $h^*$，故
\[
 a_k=\gamma C_k h^*,\qquad
 \frac{a_{k,n}}{a_{k,m}}=\frac{h_n^*}{h_m^*}
 \quad(h_m^*\ne0).
\]

令 $H=\iota^{\mathsf T}h^*$。若 $H\ne0$，共同风险基金的内部权重为 $z_R=h^*/H$，其权重之和为 $1$。参与者 $k$ 对该基金投入 $\gamma C_kH$，就能实现原来的风险头寸
\[
 (\gamma C_k H)z_R=\gamma C_k h^*=a_k.
\]

剩余的 $w_k-\gamma C_kH$ 投入无风险证券，由此得到两基金货币分离。
\end{solution}

若卖空或财富下限等额外约束绑定，需重新检查缩放后的可行集是否仍然共同。

CARA 效用也可得到共同风险投资方向。

\begin{solution}[label-content={证明},blank-type=none]
若 $u_k'(t)=\exp(-t/d_k)$、$d_k>0$，则一阶条件为
\[
 e^{-w_kx_F/d_k}\mathrm E[e^{-(a_k/d_k)^{\mathsf T}\tilde q}\tilde q]=0.
\]

所以参与者可以选取相同的 $a_k/d_k$，再按相同方法检查该方向的净投资额并归一化。
\end{solution}

\begin{solution}[label-content={证明},blank-type=none]
\textbf{用财富和权重表示同一变换。}对有限的 $\gamma\ne0$，设 $x_F>0$，定义
\[
 W_k'=w_k+\frac{\gamma d_k}{x_F},\qquad
 z_{k,n}'=\frac{w_k}{W_k'}z_{k,n}\ (n<N),\qquad
 z_{k,N}'=\frac{w_k}{W_k'}z_{k,N}+1-\frac{w_k}{W_k'}.
\]
由 $W_k'/\gamma=C_k/x_F>0$，变换后的权重有定义，且 $\sum_n z_{k,n}'=1$。将 $d_k+w_k z_k^{\mathsf T}\tilde x/\gamma=(W_k'/\gamma)(z_k')^{\mathsf T}\tilde x$ 代入\WJEqRef{eq:wj-8-13}{式~}[（以无风险证券为基准的一阶条件）]\footnote{第八讲在内点最优、头寸可双向微调且可在期望号内求导时，把一元从无风险证券移到风险证券 $n$，期末财富就增加 $\tilde x_n-x_F$。这笔转移的期望边际效用为零，即 $\mathrm E[u_k'(\tilde w_k)(\tilde x_n-x_F)]=0$。净收益率之差与总收益率之差相同。}，并除去正因子 $(W_k'/\gamma)^{-\gamma}$，得到
\[
 \mathrm E\!\left[\bigl((z_k')^{\mathsf T}\tilde x\bigr)^{-\gamma}
                    (\tilde x_n-x_F)\right]=0,
 \qquad n=1,\ldots,N-1.
\]
变换后的定义域为 $(z_k')^{\mathsf T}\tilde x>0$，方程和定义域均与 $w_k,d_k$ 无关，因此可以选择同一组 $z_k'$。再还原
\[
 z_{k,n}=\frac{W_k'}{w_k}z_{k,n}'\ (n<N),\qquad
 z_{k,N}=\frac{W_k'}{w_k}z_{k,N}'+1-\frac{W_k'}{w_k}.
\]
\end{solution}

风险证券之间的相对权重由共同风险方向决定。若共同风险方向为零，全投无风险证券即可。若方向非零而净投资额为零，多空头寸恰好抵消。可以把这一零成本头寸叠加在无风险证券上，构造另一只单位成本基金，与无风险证券一起实现一般两基金分离。\WJRef{exer:wj-nine-7}{习题~} 给出了这种情形的例子。

\textbf{必要性。}教材引用 Cass--Stiglitz（1970）。这里考察的是对所有投资机会都成立的偏好刻画；在某个给定市场观察到货币分离，不能据此反推出共同的 HARA 形状参数。使用反向结论时，须核对原文的市场、偏好与可行域条件。

\section{共同基金分离和风险投资}\label{sec:wj-nine-3}
两基金货币分离把 $N-1$ 只风险证券的问题化为\textbf{一只共同风险基金}与无风险证券之间的选择。令共同风险基金的净收益率为 $\tilde r_Q$，风险投资总额为 $a_Q$，则
\[
 \tilde w=w(1+r_F)+a_Q(\tilde r_Q-r_F),\qquad a_F=w-a_Q.
\]

以下固定共同风险基金，假设 $\mathrm E[\tilde r_Q]>r_F$，并沿用\WJRef{sec:wj-eight-single}{第~}[~节（单只风险证券）]的模型条件。此时最优风险基金投资金额为正。可以分别比较同一财富下的不同偏好，以及同一偏好下的不同财富。

\textbf{固定财富，比较两位参与者。}在相同财富下，绝对风险厌恶程度较高的人在共同风险基金上的最优投资金额不高于另一人。

\begin{solution}[label-content={证明},blank-type=none]
设参与者 1 的绝对风险厌恶系数在相关财富区间上处处不低于参与者 2。由\WJRef{thm:wj-seven-4}{定理~}[（Pratt 定理）]，$u_1=f\circ u_2$，其中 $f$ 递增且凹。让参与者 1 先持有参与者 2 的最优风险投资金额。此时两人面对相同的期末财富 $\tilde w$，且 $u_1'(\tilde w)=f'(u_2(\tilde w))u_2'(\tilde w)$。记 $b=w(1+r_F)$、$D=\tilde r_Q-r_F$、$K=f'(u_2(b))$，并以 $F_1$ 表示参与者 1 对风险投资金额的期望效用。参与者 2 在 $a_2^*$ 处满足\WJEqRef{eq:wj-8-13}{式~}，即 $\mathrm E[u_2'(\tilde w)D]=0$。从参与者 1 的导数中减去这一等式的 $K$ 倍，得到
\[
 F_1'(a_2^*)=
 \mathrm E\!\left[\bigl(f'(u_2(\tilde w))-K\bigr)u_2'(\tilde w)D\right].
\]
由于 $\tilde w=b+a_2^*D$ 且 $a_2^*>0$，$D>0$ 时 $\tilde w>b$，递增的 $u_2$ 和非递增的 $f'$ 使括号非正。$D<0$ 时括号非负。故各状态的乘积均非正，$F_1'(a_2^*)\le0$。由目标函数的凹性，参与者 1 的最优风险投资金额不高于参与者 2，无风险投资金额则不低于参与者 2。
\end{solution}

\textbf{固定偏好，比较不同财富。}由\WJRef{thm:wj-eight-3}{定理~}[（投资金额的财富效应）]与\WJRef{thm:wj-eight-4}{定理~}[（投资比例的财富效应）]\footnote{比较时固定基金收益的分布和无风险利率，要求风险溢价为正，最优解在所比较的财富范围内始终为可微的内点解，且无风险基准财富和各状态的最优期末财富均在同一正财富区间。这里用的是从风险厌恶的变化推断投资的变化；反过来推断偏好时，第八讲要求结论对所有符合条件的投资机会都成立。}，
\begin{itemize}
 \item \hyperlink{wj-risk-aversion-wealth-changes}{\textcolor{WJReferenceColor}{DARA}}、\WJLink{sec:wj-seven-cara}{CARA}、\hyperlink{wj-risk-aversion-wealth-changes}{\textcolor{WJReferenceColor}{IARA}} 分别使风险基金投资\textbf{金额}随财富增加、保持不变、减少。
 \item \hyperlink{wj-risk-aversion-wealth-changes}{\textcolor{WJReferenceColor}{DRRA}}、\WJLink{sec:wj-seven-crra}{CRRA}、\hyperlink{wj-risk-aversion-wealth-changes}{\textcolor{WJReferenceColor}{IRRA}} 分别使风险基金投资\textbf{比例}随财富增加、保持不变、减少。
\end{itemize}

无风险投资金额与比例分别由以下关系决定
\[
 \frac{da_F}{dw}=1-\frac{da_Q}{dw},\qquad
 \frac{a_F}{w}=1-\frac{a_Q}{w}.
\]

新增财富中用于风险基金的金额小于新增财富时，无风险投资金额上升。因此即使 DARA 保证 $da_Q/dw>0$，也要把这个导数与 $1$ 比较，才能判断 $a_F$ 的变化。若相对风险厌恶随财富非递减，风险基金比例非递增，无风险投资比例就非递减。一般的多风险证券模型还需同时分析风险组合内部权重的变化。

\Needspace{13\baselineskip}
\WJChapterSummary
\begin{tcolorbox}[summary,breakable=false]
\begin{enumerate}
 \item FSD 是所有递增效用都认可的排序，等价于累积分布函数处处更低；它推出均值不低，但均值不低不能反推 FSD。
 \item 本教材的 SSD 比较等均值风险，等价于累积分布函数之差的累计积分处处非正，也等价于较差分布是较好分布的均值保持扩散。
 \item 逐状态占优与随机占优须区分：前者在允许卖空时导出套利，后者本身不保证存在套利；单证券的随机占优也不能直接延伸到任意风险组合。
 \item $F$ 基金分离使每位参与者都能在 $F$ 只共同基金中找到最优组合。其中一只为无风险证券、其余基金只由风险证券组成时，称为货币分离。
 \item 收益分布的因子结构，或参与者偏好的共同结构，可以带来两基金分离。两基金货币分离使单风险资产的投资结论适用于共同风险基金。
\end{enumerate}
\end{tcolorbox}

\fi

\WJChapterExercises
% !TeX root = ../../金融经济学_王江_讲义及习题.tex
% 09_组合选择二 课后习题（由章节文件拆分）

\begin{problem}
\WJNineExercise{exer:wj-nine-1}{} 假设有四个等概率的可能状态。状态空间是 $\Omega=\{\omega_{1},\cdots,\omega_{4}\}$。考虑风险资产 $A$ 和 $B$，它们的收益率如下：
\begin{center}
\begin{tabular}{c|cccc}
    & $\omega_{1}$ & $\omega_{2}$ & $\omega_{3}$ & $\omega_{4}$ \\
    \hline
    $\tilde{r}_{A}$ & 0.5 & 0.5 & 0.7 & 0.7 \\
    $\tilde{r}_{B}$ & 0.9 & 0.8 & 0.4 & 0.3
\end{tabular}
\end{center}
在随机占优的意义上，你能为它们排序吗？
\end{problem}
\Needspace{5\baselineskip}
\begin{solution}
\[
    F_{A}(x)=\begin{cases}
        0 & x<0.5\\
        0.5 & 0.5\le x<0.7\\
        1 & 0.7\le x
    \end{cases}
\]
\[
    F_{B}(x)=\begin{cases}
        0 & x<0.3\\
        0.25 & 0.3\le x<0.4\\
        0.5 & 0.4\le x<0.8\\
        0.75 & 0.8\le x<0.9\\
        1 & 0.9\le x
    \end{cases}
\]
\[
    S(y)=\int_{-\infty}^{y}\bigl(F_{A}(x)-F_{B}(x)\bigr)\,dx=\begin{cases}
        0 & y<0.3\\
        -0.25(y-0.3) & 0.3\le y<0.4\\
        -0.025-0.5(y-0.4) & 0.4\le y<0.5\\
        -0.075 & 0.5\le y<0.7\\
        -0.075+0.5(y-0.7) & 0.7\le y<0.8\\
        -0.025+0.25(y-0.8) & 0.8\le y<0.9\\
        0 & 0.9\le y
    \end{cases}
\]

这里 $F_A,F_B$ 是净收益率的分布函数；给两个收益率同时加上 $1$，占优关系保持不变。

在 $x=0.4$ 时 $F_A<F_B$，在 $x=0.7$ 时 $F_A>F_B$，两条分布函数交叉，故两个方向的一阶随机占优均不成立。又有
\[
 \mathrm E[\tilde r_A]=\mathrm E[\tilde r_B]=0.6,\qquad S(y)\le0\quad(\forall y),
\]

由\ifdefined\WJHandoutKnowledge 定理~\ref{thm:wj-nine-3}\else 二阶随机占优判据\fi（等均值下，分布函数之差的累计积分非正可判定二阶随机占优），$A\succeq_{\mathrm{SSD}}B$。
\end{solution}

\Needspace{18\baselineskip}
\begin{problem}
\WJNineExercise{exer:wj-nine-2}{*} 两只风险资产的收益率分别为 $\tilde{r}_{A}$ 和 $\tilde{r}_{B}$。考虑一个投资者，他的初始财富为 $w=1$、效用函数为 $u(\cdot)$，其中 $u'(\cdot)>0$ 且 $u''(\cdot)<0$。他最大化期望效用 $E[u(\tilde{w})]$。令他的最优组合为 $[a,1-a]$。
\begin{enumerate}[label=(\alph*)]
    \item 假设 $\tilde{r}_{A}=\tilde{x}$，$\tilde{r}_{B}=\tilde{x}+\tilde{e}$，其中 $\tilde{e}$ 和 $\tilde{x}$ 互相独立且 $E[\tilde{e}]=0$。
    \begin{itemize}
        \item $\tilde{r}_{A}$ 是否二阶随机占优于 $\tilde{r}_{B}$？
        \item 对于投资者的最优投资组合 $a$，你能得到什么样的结论？
    \end{itemize}
    \item 假设 $\tilde{r}_{A}=\tilde{x}$，$\tilde{r}_{B}=\tilde{y}+\tilde{e}$，$\tilde{x}$、$\tilde{y}$、$\tilde{e}$ 互相独立，$\tilde{x}\sim_{d}\tilde{y}$ 且 $E[\tilde{e}]=0$。
    \begin{itemize}
        \item $\tilde{r}_{A}$ 是否二阶随机占优于 $\tilde{r}_{B}$？
        \item 证明投资者在资产 $A$ 上的投资额要大于他资产 $B$ 上的投资额，即 $a>1/2$。
    \end{itemize}
    \item 解释上面两种情形下结论之间的差异。
\end{enumerate}
\end{problem}
\Needspace{5\baselineskip}
\begin{solution}[show-qed=false]
沿用净收益率记号，期末财富为
\[
 \tilde w=1+a\tilde r_A+(1-a)\tilde r_B.
\]
\begin{enumerate}[label=(\alph*)]
 \item $\tilde r_B=\tilde x+\tilde e$ 且 $\mathrm E[\tilde e\mid\tilde x]=0$，所以 $A\succeq_{\mathrm{SSD}}B$。对任意可行的 $a$，条件 Jensen 不等式给出
 \[
 \begin{aligned}
 \mathrm E[u(\tilde w)]
 &=\mathrm E\{\mathrm E[u(1+\tilde x+(1-a)\tilde e)\mid\tilde x]\}\\
 &\le\mathrm E[u(1+\tilde x)].
 \end{aligned}
 \]

 全部投资于 $A$（$a=1$）达到上界。若 $\tilde e$ 非退化，严格凹性使 $a\ne1$ 时不等式严格，故 $a^*=1$ 唯一；若 $\tilde e=0$，两资产收益相同，任意可行配置均最优。

 \item 由独立性，$\tilde y+\tilde e$ 是 $\tilde y$ 的均值保持扩散，而 $\tilde x\sim_d\tilde y$，故仍有 $A\succeq_{\mathrm{SSD}}B$。记
 \[
 V(a)=\mathrm E[u(1+a\tilde x+(1-a)\tilde y+(1-a)\tilde e)],
 \]
 则在可微且可交换求导与期望时，
 \[
 \begin{aligned}
 V'(a)&=\mathrm E[u'(\tilde w)(\tilde x-\tilde y-\tilde e)],\\
 V''(a)&=\mathrm E[u''(\tilde w)(\tilde x-\tilde y-\tilde e)^2]\le0.
 \end{aligned}
 \]

 令 $S=\tilde x+\tilde y$。给定 $(S,\tilde e)$ 后，$\tilde x,\tilde y$ 的交换对称性给出 $\mathrm E[\tilde x-\tilde y\mid S,\tilde e]=0$。因此
 \[
 \begin{aligned}
 V'(\tfrac12)
 &=-\mathrm E[u'(1+S/2+\tilde e/2)\tilde e]\\
 &=-\mathrm E\!\left[\operatorname{Cov}
   \bigl(u'(1+S/2+\tilde e/2),\tilde e\mid S\bigr)\right]>0.
 \end{aligned}
 \]

 最后的严格不等式使用 $\tilde e$ 非退化、独立于 $S$，以及 $u'$ 严格递减。另一方面，若 $\tilde x$ 非退化，
 \[
 V'(1)=\operatorname{Cov}(u'(1+\tilde x),\tilde x)<0.
 \]

 当区间 $[0,1]$ 上的财富在效用定义域内、上述期望有限时，严格凹性给出唯一的 $a^*\in(1/2,1)$。原题的严格比较以这些非退化条件为前提；例如 $\tilde e=0$ 且 $\tilde x$ 非退化时，最优配置为 $a^*=1/2$。

 \item 第一种情形中，两资产共享完全相同的基础收益，持有 $B$ 只增加零均值噪声，因此最优选择为全部投资 $A$。第二种情形中，独立的基础收益带来分散化收益，而额外噪声降低了 $B$ 的吸引力，最优权重因而偏向 $A$。
\end{enumerate}
\end{solution}

\begin{problem}
\WJNineExercise{exer:wj-nine-3}{} 证券市场最初由 $N$ 只证券构成。证明：当把另外 $M$ 只证券添加到市场中去时，如果投资者的最优投资组合中包含新证券，那么他的福利（即他的期望效用）会增加（也就是说，更多选择总是好的）。
\end{problem}
\Needspace{5\baselineskip}
\begin{solution}
设新证券加入前后的可行组合集合分别为 $\mathcal A_0\subseteq\mathcal A_1$，则
\[
 \max_{a\in\mathcal A_1}\mathrm E[u(\tilde w(a))]
 \ge\max_{a\in\mathcal A_0}\mathrm E[u(\tilde w(a))].
\]

可行集扩大使最优福利至少保持原水平。只有原有组合都达不到新可行集中的最优福利时，福利才严格提高。

例如，新证券与原有证券的价格和支付完全相同。投资者可以把原持仓的一部分换成新证券，期末财富与期望效用保持不变，尽管新组合中持有新证券。

若新市场中最优的期末财富唯一，并且旧证券组合无法实现这一财富，则福利严格提高。否则，旧市场最优组合也会在新市场最优，与上述唯一性及不可实现性矛盾。练习~\ref{exer:wj-nine-8} 在正定协方差条件下给出严格改善的具体计算。
\end{solution}

\begin{problem}
\WJNineExercise{exer:wj-nine-4}{} 假设有 $N$ 只证券，它们的供给都是无限的、收益率为 $\tilde{r}_{i}$，$i=1,\cdots,N$。经济中的参与者都是不满足的、风险中性的且具有正的禀赋。市场不允许卖空。在什么条件下 $K$ 基金分离成立（$K<N$）？
\end{problem}
\Needspace{5\baselineskip}
\begin{solution}
假设参与者采用共同概率信念（原题未明写此条件），记 $\bar r_{\max}=\max_i\mathrm E[\tilde r_i]$，并令
\[
 I=\{i:\mathrm E[\tilde r_i]=\bar r_{\max}\},\qquad m=|I|.
\]

风险中性且禁止卖空时，最优权重集合为
\[
 \mathcal Z^*=\left\{z:z_i\ge0,\ \sum_i z_i=1,\ z_i=0\ (i\notin I)\right\}.
\]

按\ifdefined\WJHandoutKnowledge 定义~\ref{def:wj-nine-4}\else 共同基金分离的定义\fi（每人均可选择一个由共同基金构成的最优组合），只需为每人选择一个共同基金中的最优组合。任选 $j\in I$，所有人全部投资于证券 $j$ 即可，故单基金分离成立，也可用包含它的 $K$ 只基金实现分离。

若要求表示\textbf{全部}最优组合，则 $\mathcal Z^*$ 是由 $m$ 个最优证券顶点生成的 $m-1$ 维单纯形，最少需要 $m$ 只基金。因此用至多 $K$ 只基金覆盖全部最优组合的条件是 $m\le K$；最少基金数恰为 $K$ 的条件是 $m=K$。
\end{solution}

\begin{problem}
\WJNineExercise{exer:wj-nine-5}{} 假设在市场上有 $N$ 只证券，它们的收益率是独立同分布的。证明单基金分离成立。
\end{problem}
\Needspace{5\baselineskip}
\begin{solution}
假设参与者效用凹、采用共同概率，且可行集合对证券置换和取平均封闭。对任意可行权重 $z$，将其全部 $N!$ 个置换取平均，得到
\[
 \frac1{N!}\sum_{\pi}z_{\pi}=\frac{\iota}{N}.
\]
独立同分布使各置换组合的收益分布相同；由 Jensen 不等式，任意可行权重组合的期望效用均不超过等权组合：
\[
 \mathrm E\!\left[u\!\left(w\left(1+\sum_{i=1}^N z_i\tilde r_i\right)\right)\right]
 \le\mathrm E\!\left[u\!\left(w\left(1+\frac1N\sum_{i=1}^N\tilde r_i\right)\right)\right].
\]

因此，每位参与者都可以选择同一只等权基金 $z_f=\iota/N$，单基金分离成立。若效用严格凹，且不同可行权重对应的期末财富以正概率不同，则等权组合还是唯一最优解。
\end{solution}

\begin{samepage}
\begin{problem}
\WJNineExercise{exer:wj-nine-6}{*} 证券市场由一只无风险债券和 $N$ 只风险证券组成。无风险证券的收益率为 $r_{F}$ 而风险证券的收益率向量为 $\tilde{r}=\begin{bmatrix}\tilde{r}_{1} \\ \ldots \\ \tilde{r}_{N}\end{bmatrix}$。另外，$E[\tilde{r}]=\bar{r}=\begin{bmatrix}\bar{r}_{1} \\ \ldots \\ \bar{r}_{N}\end{bmatrix}$ 且 $E[(\tilde{r}-\bar{r})(\tilde{r}-\bar{r})^{\top}]=\Sigma$（$\Sigma$ 是风险证券收益率的协方差阵）。考虑一个财富为 $w$ 的参与者的投资问题。假设他具有二次效用函数 $u(\tilde{w})=E[\tilde{w}]-\tfrac12 a E[\tilde{w}^{2}]$。
\begin{enumerate}[label=(\alph*)]
    \item 求解他的最优投资组合。
    \item 证明对于具有二次效用函数的参与者，两基金货币分离成立。
\end{enumerate}
\end{problem}
\end{samepage}
\Needspace{5\baselineskip}
\begin{solution}
原题的 $E[\tilde w]-aE[\tilde w^2]/2$ 表示期望效用，其财富效用为 $u(t)=t-at^2/2$。设 $w>0$、$a>0$，且 $\Sigma$ 正定；交易允许自由借贷和卖空。二次效用递增区间为 $t<1/a$，下式求解无额外财富约束的二次目标。
\begin{enumerate}[label=(\alph*)]
    \item 令风险资产权重为 $x$，并把期末财富写成与练习~\ref{exer:wj-eight-2} 相同的形式
    \[
        \tilde{w}=w\bigl(1+r_{F}+x^{\top}(\tilde{r}-r_{F}\iota)\bigr)
    \]

    记 $\mu=\bar{r}-r_{F}\iota$。则 $E[\tilde{w}]=w(1+r_{F}+x^{\top}\mu)$，且
    \[
        E[\tilde{w}^{2}]=w^{2}\bigl[(1+r_{F}+x^{\top}\mu)^{2}+x^{\top}\Sigma x\bigr]
    \]

    对 $x$ 求导，一阶条件为
    \[
        w\mu-aw^{2}\bigl[(1+r_{F}+x^{\top}\mu)\mu+\Sigma x\bigr]=0
    \]
    目标函数的 Hessian 为 $-aw^2(\mu\mu^{\mathsf T}+\Sigma)\prec0$，故一阶条件给出唯一最大值。即 $\bigl[1-aw(1+r_{F})\bigr]\mu=aw(\mu\mu^{\top}+\Sigma)x$，解得最优权重
    \[
        x^{*}=\frac{1-aw(1+r_{F})}{aw}\bigl[(\bar{r}-r_{F}\iota)(\bar{r}-r_{F}\iota)^{\top}+\Sigma\bigr]^{-1}(\bar{r}-r_{F}\iota)
    \]
    其中，$\iota$ 是元素全为 1 的 $N$ 维列向量。风险组合中各资产的相对权重由 $(\mu\mu^{\top}+\Sigma)^{-1}\mu$ 决定，与财富、风险厌恶的水平只差一个共同标量。

    \item 记 $d=(\mu\mu^{\mathsf T}+\Sigma)^{-1}\mu$，所有人的风险头寸均为 $d$ 的标量倍数。若 $D=\iota^{\mathsf T}d\ne0$，定义仅由风险证券构成的基金 $z_R=d/D$，每人便在 $z_R$ 与无风险证券之间配置，得到两基金货币分离。

    若 $d=0$，最优选择为全部投资于无风险证券。若 $d\ne0$ 而 $D=0$，共同方向是零成本风险头寸，可用无风险证券 $F$ 与单位成本基金 $F+(d,0)$ 实现一般两基金分离；此时风险头寸本身的净投资额为零，无法按 $d/D$ 归一化为纯风险基金。
\end{enumerate}
\end{solution}

\Needspace{11\baselineskip}
\begin{problem}
\WJNineExercise{exer:wj-nine-7}{*} 考虑练习~\ref{exer:wj-nine-6} 中描述的市场。另外，我们假设风险收益率是联合正态分布的且参与者的绝对风险厌恶系数恒为 $a$。
\begin{enumerate}[label=(\alph*)]
    \item 求解他的最优投资组合。
    \item 证明对于具有常数绝对风险厌恶的参与者，两基金货币分离成立。
\end{enumerate}
\end{problem}
\Needspace{5\baselineskip}
\begin{solution}
设 $w>0$、$a>0$，且 $\Sigma$ 正定；交易允许自由借贷和卖空。
\begin{enumerate}[label=(\alph*)]
    \item 仍令风险资产权重为 $x$，
    \[
        \tilde{w}=w\bigl(1+r_{F}+x^{\top}(\tilde{r}-r_{F}\iota)\bigr)
    \]
    $\tilde{r}$ 联合正态则 $\tilde{w}$ 正态。与练习~\ref{exer:wj-eight-3}（指数效用下正态财富的期望效用计算）相同，
    \[
        E\bigl[-e^{-a\tilde{w}}\bigr]=-\exp\Bigl(-aE[\tilde{w}]+\tfrac12 a^{2}\operatorname{Var}(\tilde{w})\Bigr)
    \]
    其中 $E[\tilde{w}]=w(1+r_{F}+x^{\top}\mu)$，$\operatorname{Var}(\tilde{w})=w^{2}x^{\top}\Sigma x$，$\mu=\bar{r}-r_{F}\iota$。最大化期望效用等价于最大化确定性等价
    \[
      CE=w(1+r_F)+wx^{\mathsf T}\mu-\frac a2w^2x^{\mathsf T}\Sigma x.
    \]
    其 Hessian 为 $-aw^2\Sigma\prec0$，因此最优解唯一。对 $x$ 的一阶条件为
    \[
        w\mu-aw^{2}\Sigma x=0
    \]

    因而最优权重为
    \[
        x^{*}=\frac{1}{aw}\Sigma^{-1}(\bar{r}-r_{F}\iota)
    \]
    其中，$\iota$ 是元素全为 1 的 $N$ 维列向量。相对权重由 $\Sigma^{-1}\mu$ 决定，与 $a$、$w$ 无关。

    \item 记 $d=\Sigma^{-1}\mu$，所有人的风险头寸均为 $d$ 的标量倍数。若 $D=\iota^{\mathsf T}d\ne0$，定义仅由风险证券构成的基金 $z_R=d/D$，每人便在 $z_R$ 与无风险证券之间配置，得到两基金货币分离。

    若 $d=0$，最优选择为全部投资于无风险证券。若 $d\ne0$ 而 $D=0$，共同方向是零成本风险头寸，可用无风险证券 $F$ 与单位成本基金 $F+(d,0)$ 实现一般两基金分离；此时风险头寸本身的净投资额为零，无法按 $d/D$ 归一化为纯风险基金。
\end{enumerate}
\end{solution}

\begin{problem}
\WJNineExercise{exer:wj-nine-8}{*} 继续练习~\ref{exer:wj-nine-7} 中的问题。
\begin{enumerate}[label=(\alph*)]
    \item 计算在最优投资组合下参与者的期望效用。
    \item 假设现在有 $M$ 支新风险证券加入到证券市场中。（现有 $N$ 只证券的收益率不能生成新证券的收益率。）所有风险证券的收益率都是联合正态分布的。证明，如果参与者的最优投资组合中含有新的证券，无论 $M$ 只新证券的收益率分布如何，参与者的福利总是增加。
    \item 当引入新证券后，计算他的福利增加的确定性等价值。
\end{enumerate}
\end{problem}
\Needspace{5\baselineskip}
\begin{solution}
设 $w>0$、$a>0$，并假设加入新证券后的风险收益协方差矩阵正定。记各风险证券的超额收益均值为 $\mu$。
\begin{enumerate}[label=(\alph*)]
 \item 由练习~\ref{exer:wj-nine-7}（指数效用与联合正态收益下的最优风险资产权重），$x^*=\Sigma^{-1}\mu/(aw)$。代回指数效用，得
 \[
 \begin{aligned}
 \mathrm E[u(\tilde w^*)]
 &=-\exp\!\left(-aw(1+r_F)-aw{x^*}^{\mathsf T}\mu
               +\tfrac12a^2w^2{x^*}^{\mathsf T}\Sigma x^*\right)\\
 &=-\exp\!\left(-aw(1+r_F)-\tfrac12\mu^{\mathsf T}\Sigma^{-1}\mu\right).
 \end{aligned}
 \]

 \item 将均值和协方差按旧、新证券分块：
 \[
 \mu=\begin{bmatrix}\mu_0\\\mu_1\end{bmatrix},\qquad
 \Sigma=\begin{bmatrix}A&C\\C^{\mathsf T}&B\end{bmatrix}.
 \]

 原题所述新证券不能由旧证券生成，未必排除新证券之间冗余；这里的扩充协方差正定是更强的条件。定义新收益剔除旧收益线性成分后的协方差及超额收益均值：
 \[
 S=B-C^{\mathsf T}A^{-1}C\succ0,\qquad
 \delta=\mu_1-C^{\mathsf T}A^{-1}\mu_0.
 \]

 由分块一阶方程，先消去旧证券权重，得到 $awSx_1^*=\delta$。分块求逆进一步给出
 \[
 \mu^{\mathsf T}\Sigma^{-1}\mu
 =\mu_0^{\mathsf T}A^{-1}\mu_0+\delta^{\mathsf T}S^{-1}\delta,
 \qquad x_1^*=\frac1{aw}S^{-1}\delta.
 \]

 新证券的最优权重向量非零，当且仅当 $\delta\ne0$；此时 $\delta^{\mathsf T}S^{-1}\delta>0$，期望效用严格提高。这个结果把新证券提供的收益机会与它们在旧证券之外增加的风险联系起来。

 \item 由 $-e^{-aCE}=\mathrm E[-e^{-a\tilde w^*}]$，最优财富的确定性等价为
 \[
 CE=w(1+r_F)+\frac1{2a}\mu^{\mathsf T}\Sigma^{-1}\mu.
 \]

 因而福利增加的确定性等价值为
 \[
 \boxed{\Delta CE=\frac1{2a}
 \left(\mu_{\mathrm{new}}^{\mathsf T}\Sigma_{\mathrm{new}}^{-1}\mu_{\mathrm{new}}
       -\mu_{\mathrm{old}}^{\mathsf T}\Sigma_{\mathrm{old}}^{-1}\mu_{\mathrm{old}}\right)
 =\frac1{2a}\delta^{\mathsf T}S^{-1}\delta.}
 \]
\end{enumerate}
\end{solution}

\begin{problem}
\WJNineExercise{exer:wj-nine-9}{*} 在练习~\ref{exer:wj-nine-7} 中，假设只有一只新证券（即 $M=1$），它的收益率是 $\tilde{r}_{N+1}$。令 $\bar{r}_{N+1}$ 为它的预期收益率、$\sigma_{N+1}$ 为收益率波动率而 $\sigma_{\cdot,N+1}$ 为它与 $N$ 只原有证券之间的协方差向量。另外，证券 $N+1$ 上的最优组合权重记作 $z_{N+1}$。
\begin{enumerate}[label=(\alph*)]
    \item 假设 $\bar{r}_{N+1}-r_{F}>(<)\,0$。我们能得到结论 $z_{N+1}>(<)\,0$ 吗？给出解释。
    \item 假设 $\bar{r}_{N+1}<\bar{r}_{n}$ 且 $\sigma_{N+1}>\sigma_{n}$，$\forall\,n=1,\ldots,N$。也就是说，其它所有证券在收益率和方差上都优于证券 $N+1$。这是否意味着 $z_{N+1}=0$ 即最优组合不应含有证券 $N+1$？给出解释。
\end{enumerate}
\end{problem}
\Needspace{5\baselineskip}
\begin{solution}
记 $A$ 为原有风险证券的协方差矩阵，$c=\sigma_{\cdot,N+1}$，$\mu_0$ 为原有证券的超额收益均值，$\mu_1=\bar r_{N+1}-r_F$。假设扩充后的协方差矩阵正定，$a,w>0$。
\begin{enumerate}[label=(\alph*)]
 \item 由练习~\ref{exer:wj-nine-8} 的分块公式（剔除旧资产线性成分后，新资产最优权重的表达式），
 \[
 \boxed{z_{N+1}=\frac{\mu_1-c^{\mathsf T}A^{-1}\mu_0}
 {aw\left(\sigma_{N+1}^2-c^{\mathsf T}A^{-1}c\right)}.}
 \]

 分母为正，持仓方向由剔除旧资产风险敞口后的超额收益
 $\mu_1-c^{\mathsf T}A^{-1}\mu_0$ 决定。因此仅凭 $\mu_1$ 的正负，无法确定 $z_{N+1}$ 的正负。例如一只旧资产取 $\mu_0=0.10$、方差 $0.04$，新资产取 $\mu_1=0.05$、方差 $0.09$，协方差 $0.04$。协方差矩阵行列式为 $0.002>0$，但残差超额收益为 $0.05-(0.04/0.04)0.10=-0.05$，故新资产最优权重为负。

 \item 均值和标准差的单独比较没有包含协方差信息。只要
 $\mu_1-c^{\mathsf T}A^{-1}\mu_0\ne0$，新证券就有非零最优持仓。

 例如取一只旧风险证券，令 $r_F=0$，旧证券的均值为 $0.10$、标准差为 $0.20$，新证券的均值为 $0.05$、标准差为 $0.30$，两者独立且联合正态。新证券的均值更低、标准差更高，但
 \[
 z_{N+1}=\frac{0.05}{aw\cdot0.09}>0.
 \]

 它提供了独立的正超额收益来源，持有它有助于分散风险。
\end{enumerate}
\end{solution}


\endgroup
